\section{Gabarito e resoluções comentadas — Noções de Direitos Humanos}

\begin{solucao}{09.01}\textbf{Errado.} Há distinção didática entre planos internacional e constitucional interno, mas eles dialogam intensamente por tratados, jurisprudência e interpretação constitucional.\end{solucao}
\begin{solucao}{09.02}\textbf{Certo.} As dimensões são cumulativas e interdependentes; não há substituição histórica dos direitos anteriores.\end{solucao}
\begin{solucao}{09.03}\textbf{Gabarito: A.} A DUDH foi proclamada em 10 de dezembro de 1948 pela Assembleia Geral da ONU.\end{solucao}
\begin{solucao}{09.04}\textbf{Certo.} A DUDH reúne direitos civis, políticos, econômicos, sociais e culturais.\end{solucao}
\begin{solucao}{09.05}\textbf{Certo.} A Agenda 2030 possui 17 ODS e 169 metas.\end{solucao}
\begin{solucao}{09.06}\textbf{Gabarito: B.} A LBI adota conceito social/relacional: impedimento de longo prazo em interação com barreiras que restringem participação.\end{solucao}
\begin{solucao}{09.07}\textbf{Certo.} Deficiência não produz incapacidade civil automática.\end{solucao}
\begin{solucao}{09.08}\textbf{Certo.} A própria LBI inclui a recusa de adaptação razoável no conceito de discriminação em razão da deficiência.\end{solucao}
\begin{solucao}{09.09}\textbf{Gabarito: A.} A Lei nº 10.098/2000 disciplina promoção da acessibilidade.\end{solucao}
\begin{solucao}{09.10}\textbf{Certo.} O Estatuto reconhece ações afirmativas como instrumentos de promoção da igualdade de oportunidades.\end{solucao}

\begin{solucao}{09.11}\textbf{Gabarito: B.} Acessibilidade alcança informação, comunicação, sistemas e tecnologias, não apenas edificações.\end{solucao}
\begin{solucao}{09.12}\textbf{Certo.} A avaliação biopsicossocial considera corpo, ambiente, aspectos pessoais, atividade e participação.\end{solucao}
\begin{solucao}{09.13}\textbf{Errado.} Igualdade material pode exigir diferenciação legítima e ajustes para neutralizar barreiras reais.\end{solucao}
\begin{solucao}{09.14}\textbf{Gabarito: B.} Adaptação razoável é ajuste necessário e adequado que não imponha ônus desproporcional ou indevido.\end{solucao}
\begin{solucao}{09.15}\textbf{Errado.} A LBI protege autonomia e plena capacidade civil, inclusive em direitos existenciais.\end{solucao}
\begin{solucao}{09.16}\textbf{Gabarito: B.} O rol vigente foi ampliado e inclui TEA, pessoas com mobilidade reduzida e doadores de sangue, entre outros.\end{solucao}
\begin{solucao}{09.17}\textbf{Errado.} Doadores de sangue têm prioridade após os demais beneficiários do rol legal e mediante comprovante com validade prevista em lei.\end{solucao}
\begin{solucao}{09.18}\textbf{Certo.} Os ODS são universais, ao contrário de leituras que os restringem a países em desenvolvimento.\end{solucao}
\begin{solucao}{09.19}\textbf{Gabarito: D.} ODS 16 trata de paz, justiça e instituições eficazes, responsáveis e inclusivas.\end{solucao}
\begin{solucao}{09.20}\textbf{Certo.} A Agenda define os ODS como integrados e indivisíveis.\end{solucao}
\begin{solucao}{09.21}\textbf{Gabarito: B.} Essa é a definição legal adotada pelo Estatuto da Igualdade Racial.\end{solucao}
\begin{solucao}{09.22}\textbf{Certo.} Ações afirmativas buscam igualdade material e correção de desigualdades persistentes.\end{solucao}
\begin{solucao}{09.23}\textbf{Certo.} Barreiras podem ser atitudinais, comunicacionais e tecnológicas mesmo quando a arquitetura é acessível.\end{solucao}
\begin{solucao}{09.24}\textbf{Gabarito: A.} Vídeos sem recursos acessíveis e formulários incompatíveis com tecnologia assistiva criam barreiras de comunicação/informação e tecnologia.\end{solucao}
\begin{solucao}{09.25}\textbf{Certo.} Desenvolvimento sustentável na Agenda 2030 integra dimensões econômica, social e ambiental.\end{solucao}

\begin{solucao}{09.26}\textbf{Errado.} Universalidade não significa absolutismo. Direitos podem ser harmonizados e sofrer restrições legítimas sob critérios constitucionais.\end{solucao}
\begin{solucao}{09.27}\textbf{Gabarito: A.} Exigir curatela de toda pessoa com deficiência reproduz justamente a presunção de incapacidade rejeitada pela LBI.\end{solucao}
\begin{solucao}{09.28}\textbf{Certo.} Curatela é medida excepcional e proporcional; não apaga personalidade nem direitos existenciais.\end{solucao}
\begin{solucao}{09.29}\textbf{Certo.} Esses testes fornecem evidências técnicas de acessibilidade digital efetiva.\end{solucao}
\begin{solucao}{09.30}\textbf{Gabarito: B.} Igualdade material exige remover a barreira e oferecer forma equivalente de acesso.\end{solucao}
\begin{solucao}{09.31}\textbf{Certo.} A recusa de tecnologia assistiva pode prejudicar exercício de direitos e integrar discriminação.\end{solucao}
\begin{solucao}{09.32}\textbf{Gabarito: C.} A DUDH parte da dignidade e igualdade e abrange amplo conjunto de direitos.\end{solucao}
\begin{solucao}{09.33}\textbf{Certo.} Esses direitos integram o catálogo da DUDH.\end{solucao}
\begin{solucao}{09.34}\textbf{Certo.} Transparência, informação e instituições responsáveis são temas diretamente associados ao ODS 16.\end{solucao}
\begin{solucao}{09.35}\textbf{Gabarito: A.} Sem desagregação não se mede diferença de acesso ou resultado entre grupos raciais.\end{solucao}
\begin{solucao}{09.36}\textbf{Errado.} Igualdade material admite e pode exigir políticas específicas, inclusive ações afirmativas.\end{solucao}
\begin{solucao}{09.37}\textbf{Certo.} O Estatuto possui cobertura temática ampla.\end{solucao}
\begin{solucao}{09.38}\textbf{Gabarito: A.} Prioridade precisa ser efetiva e compatível com autonomia e apoio necessário ao titular.\end{solucao}
\begin{solucao}{09.39}\textbf{Certo.} O símbolo é opcional; sua ausência não elimina direitos.\end{solucao}
\begin{solucao}{09.40}\textbf{Certo.} Alterações recentes incorporaram necessidades complexas de comunicação e recursos aumentativos/alternativos em contextos específicos.\end{solucao}

\begin{solucao}{09.41}\textbf{Gabarito: A.} O problema começa no planejamento/requisitos e se materializa em exclusão digital e possível discriminação na prestação do serviço.\end{solucao}
\begin{solucao}{09.42}\textbf{Certo.} Totais agregados podem mascarar quem efetivamente acessa o programa.\end{solucao}
\begin{solucao}{09.43}\textbf{Gabarito: A.} A demora cria acesso materialmente inferior para determinado grupo e precisa ser justificada e corrigida.\end{solucao}
\begin{solucao}{09.44}\textbf{Certo.} Acessibilidade é condição de uso com segurança e autonomia, não mera presença nominal de recurso.\end{solucao}
\begin{solucao}{09.45}\textbf{Gabarito: B.} Vincular um programa nominalmente a um ODS não demonstra causalidade, desempenho ou impacto.\end{solucao}
\begin{solucao}{09.46}\textbf{Certo.} A Agenda orienta políticas e avaliação, sem transformar automaticamente cada meta em pretensão judicial autoexecutável.\end{solucao}
\begin{solucao}{09.47}\textbf{Gabarito: B.} A redação vigente atribui aos doadores prioridade após os demais beneficiários do rol.\end{solucao}
\begin{solucao}{09.48}\textbf{Certo.} Acessibilidade é multidimensional.\end{solucao}
\begin{solucao}{09.49}\textbf{Gabarito: B.} Igualdade material exige examinar efeitos concretos e possíveis mecanismos de exclusão indireta.\end{solucao}
\begin{solucao}{09.50}\textbf{Certo.} Auditoria moderna de políticas pode integrar conformidade, equidade, acessibilidade e efetividade.\end{solucao}
